% !TeX root = ../Masterarbeit.tex

\documentclass[Masterarbeit.tex]{subfile}

\begin{document}
	
	\chapter{Data and Methods} \label{data_methods}
	
	\section{Data} \label{data}
	
	In order to gain a new estimate of North Atlantic ocean heat content, this work will use two different sources of training and testing data for the utilized neural network. On the one hand, data assimilation experiments \citep{brune_preserving_2020} from the Max Planck Institute Earth System Model (MPI-ESM, low resolution, \cite{mauritsen_developments_2019}) are deployed to train the neural network to reconstruct temperature maps. On the other hand, observational data from the fourth version of the Met Office Hadley Center 'EN' series (EN4, \cite{good_en4_2013}) subsurface ocean temperature profiles are applied as basis for the final reconstruction. After the neural network reconstructions, the temperature profiles are converted into a two dimensional OHC map via integration:
	
	
	\begin{align}
	H = c_{p} \int_{z_0}^{z_1} \rho_z \cd T_z dz ,
	\end{align}
	
	In this case $H$ denotes the ocean heat content, $c_{p}$ the specific heat capacity of sea water, $\rho_z$ the density of the water at depth $z$ and $T_z$ the temperature at this depth. This OHC field is then reduced to one value by summing over all relevant gridpoints for the calculation of an OHC timeseries.	
	
	
	\subsection{Data Assimilation} \label{data:assim}
	
	In general, Kalman Filter (KF) assimilations assume the uncertainties of model and observations to be Gaussian and then assimilates the observations sequentially. Both, model and observations, are described as probability density functions (PDFs) and the Kalman filter alternates between forecasting and analyzing the PDF's covariance and mean \cite{carrassi_data_2018}. By introducing an ensemble of model trajectories instead of one single representation of the model dynamics, the KF can be adjusted and extended in what is called an ensemble Kalman filter (EnKF). This allows for the covariance matrices to be approximated by an ensemble of model states. The covariance matrix and ensemble mean are then sequentially forecast and updated.\\
	
	For this work, an ensemble coupled MPI-ESM assimilation \citep{brune_mpi-esm-lr_1201p5_2021, brune_preserving_2020} is used to train the U-net to reconstruct masked temperature profiles. An oceanic localized EnKF, implemented with the Parallel Data Assimilation Framework \citep{nerger_software_2013}, assimilates monthly EN4 temperature and salinity profiles. Simultaneously, atmospheric nudging of ERA40/ERAInterim/ERA5 reanalyses \citep{uppala_era-40_2005, dee_era-interim_2011, hersbach_era5_2020} takes place. Monthly mean temperatures are taken from the assimilation from January 1958 -- October 2020 down to a depth of 2.000 m in 40 depth layers.\\
	
	\subsection{Observational Data} \label{data:en4}
	
	The EN4 project describes a collection of global ocean temperature and salinity profile datasets spanning from 1900 until present day. It contains observational data from Argo floats, Arctic Synoptic Basinide Oceanography (ASBO), the Global Temperature and Salinity Profile Program (GTSPP) and the World Ocean Database (WOD13) \citep{good_en4_2013}. As these datasets partly overlap, input data is preprocessed and duplicates sorted out before the complete EN4 dataset is created. The same procedure is applied to profiles that are not necessarily identified to be the same but that are very close to each other in time and space \citep{good_en4_2013}. Additionally, various quality control checks are performed in order to identify data profiles with faulty location or measurement data, such as missing depth steps or the location assignment to a point on land rather than in the ocean. The EN4 temperature profiles possess the same vertical and horizontal resolution as the assimilation data. All datapoints are brought together in NetCDF files and are publicly available under \url{https://www.metoffice.gov.uk/hadobs/} \citep{good_en4_2013}. 
	
	\subsection{Other OHC Estimates} \label{data:other_estimates}
	
	\begin{figure} [h]
		\centering
		\includegraphics[width=1.0\linewidth]{plots/all_estimates}
		\caption{Comparison of EnKF assimilation ensemble mean (red) with IAP (orange) and EN4 (purple) SPG OHC estimates. The shaded area around the assimilation ensemble mean represents its ensemble spread.}
		\label{fig:allestimates}
	\end{figure}
	
	
	In order to contextualize the reconstruction results outside of the EnKF assimilation, two other OHC estimates will be taken as reference points. Figure \ref{fig:allestimates} shows these in comparison to the EnKF assimilation ensemble mean and its spread. The EN4 project provides an interpolated version of the observed data points in the form of an objective analysis. This analysis is created by combining a persistence based forecast from the previous month's analysis with the observed temperature profiles \cite{good_en4_2013}. As the persistence relaxes towards the baseline climatology when not corrected by observations, the EN4 reanalysis tends to posses a lot of uncertainty in times with low observational density. The applied baseline climatology is taken from the years 1970 -- 2000. Thus, the overall cooling trend in the NA during that time (\ref{fig:allestimates}) leads to a cool bias in the EN4 OHC estimate especially in the early pre-Argo era \cite{good_en4_2013}. Despite these shortcomings, the EN4 analysis will be taken as a reference point in the evaluation of the network reconstructions as it is highly consistent with the observational data, especially in the well observed Argo era. Therefore, it can help provide additional context in particular to the direct neural network reconstruction of EN4 observations.\\
	In order to better contextualize the reconstruction results, the Institute for Atmospheric Physic's (IAP) OHC estimate is additionally provided as a point of comparison. The IAP mainly uses the same sources for its observational input as the EN4 objective analysis \citep{cheng_iap_2016}. The main difference to the EN4 reanalysis consists in the mapping of the unobserved gridpoints. In opposition to the statistical mapping of the EN4, which is only supplemented by an estimation of covariances and a relaxation term towards the baseline climatology, the IAP also incorporates model data to reconstruct missing data points. The IAP combines CMIP5 multi-model simulations with ensemble interpolation methods, which produces a data assimilation that is less biased even in regions with low observations compared to simpler interpolations \cite{cheng_iap_2016}. According to \cite{cheng_iap_2016}, the IAP OHC estimate provides reliable results without a significant bias from 1955 on. However, it only agrees with estimates of top-of-atmosphere radiative imbalance since 1985. Therefore, all results before 1985, but especially before 1955 possess larger uncertainties. All IAP temperature fields are publicly available and have been downloaded from \url{http://www.ocean.iap.ac.cn/pages/dataService/dataService.html?languageType=en&navAnchor=dataService}. \\
	
	
	\section{Methods} \label{methods}
	
	\subsection{The Convolutional U-net} \label{method:unet}
	
	The U-net adopted in this work is based on the one used by \cite{kadow_artificial_2020} to reconstruct near-surface air temperature data. \cite{kadow_artificial_2020} adopted image inpainting technology developed by \cite{liu_image_2018}. This original U-net focused on infilling irregular holes in a variety of images. The concept of infilling irregular holes aims at producing smooth results consistent with the holes' environment without limiting itself to rectangular holes. \cite{liu_image_2018} achieved that by designing a neural network, which utilizes partial convolutions. Partial convolutions consist of a masked and re-normalized convolution operation, which is then followed by an automatic mask-update step \citep{liu_image_2018}. In mathematical terms, this means that the partial convolution at each location is defined as
	
	\begin{align}
	x' = \begin{cases}
	W^{T} (X \cdot M) \frac{\sum(1)}{\sum(M)} + b, &\text{if} \, \sum(M) > 0\\
	0, &\text{otherwise ,}\\
	\end{cases}
	\end{align}
	
	where $W$ denotes the convolutional filter weights, $b$ its corresponding bias, $X$ the feature values for the current convolution, $M$ the corresponding mask and $1$ an array with the same size as $M$ filled with ones \cite{liu_image_2018}. This mask is then updated for each element according to \cite{liu_image_2018} as
	
	\begin{align}
	m' = \begin{cases}
	1, & \text{if} \sum(M) > 0 \\
	0, & \text{otherwise .} \\
	\end{cases}
	\end{align}
	
	With enough partial convolutions, the mask is eventually updated to incorporate only ones and the output image does not possess any empty grid points anymore. \cite{liu_image_2018} incorporated partial convolutions into the bigger framework of convolutional U-nets. U-nets downscale their input images with convolution and pooling operations in order to filter out the import data contained in its input in a so called encoding path. Afterwards, the downsampled data is brought back to its initial resolution through interpolation and deconvolution operations in a decoding path. This setup, where data is first condensed into a sort of bottleneck and upscaled again afterwards, is depicted in figure \ref{fig:unet}. The name U-net originates from its characteristical shape.
	
	\begin{figure}
		\centering
		\includegraphics[width=0.8\linewidth]{plots/Unet_infilling}
		\caption{Depiction of the original U-net utilized by \cite{liu_image_2018} for the inpainting of irregular holes. In the case of this thesis, the masked images of photographs are replaced by masked assimilation data or EN4 observations.}
		\label{fig:unet}
	\end{figure}
	
	
	In each encoding and decoding layer of the U-net a partial convolution operation is performed. This approach yields significant upsides compared to normal convolutional layers. Convolutional U-nets rely very heavily on the initial values given to the U-net in the places of the image holes. This often makes computationally expensive postprocessing necessary to compensate for the artifacts stemming from this initial value bias \citep{liu_image_2018}. Partial convolutional layers account for the initial value problem by using the described masked convolution operation followed by the automatic mask update step. This results in the reconstruction being only dependent on the input data that is not masked out and thus renders the initial hole value irrelevant \citep{liu_image_2018}. An additional upside of this approach is the fact that this dependence on only the valid input data renders the shape of the holes irrelevant. Both problems are also applicable to the reconstruction of missing temperature observations, as the missing data usually differs at every timestep in its shape. Thus, these partial convolutional U-nets can learn spatial patterns of temperature data and apply them to highly varying observational patterns \cite{kadow_artificial_2020}.\\
	Each layer of the U-net is activated using a Rectified Linear Unit (ReLU) activation functions. ReLU activation functions provide an easy possibility to reduce the amount of active neurons in each layer, which significantly speeds up neural networks compared to sigmoid or tanh activation functions. The U-net's training loss is calculated at each iteration by separating the overall loss function into five parts, each for a different feature of the network output. The most impactful factors in connection to the U-net's overall performance when reconstructing missing climate information  are the gridpoint wise loss functions for the infilled holes and originally unmasked data points \citep{plesiat_personal_2022}. The individual components are always calculated as mean absolute error (MAE) between the network output and the ground truth of the training data. The weights of the individual loss components can be adjusted and are determined manually by tryouts without any specific hyperparameter training. \\
	
	\subsection{Neural Network Setup and Parameters}
	
	The U-net conceptually described in section \ref{method:unet} can be setup in a number of different ways depending on the exact purpose of the training and the input data. In contrast to \cite{kadow_artificial_2020}, this work focuses on three dimensional temperature profiles instead of just SSTs. Thus, instead of one network, there are several at work, with each one training to reconstruct an individual depth layer. After each layer is reconstructed, the resulting filled in temperature slices can be reassembled and postprocessed. As the aim is to calculate the upper 700 m ocean heat content, the first 20 depth layers of the assimilation training data are selected, which represent temperature fields up to a depth of 700 m. \\
	The hyperparameters for each U-net, including network depth and loss parameters can be adjusted individually previous to each training instance. The default setup consists of four encoding and four pooling layers because until now, a further deepening of the neural network slows down calculations without providing significantly better results. Furthermore, the time period as well as the spatial extension in the form of longitudes and latitudes of the input data can be individually determined for each training process. In regard to the time period of the training data, two different arguments can be adjusted. On the one hand, the origin of the data itself can be adjusted, with default options from the Argo era, pre-Argo era or the entire assimilations time period. On the other hand, the time of origin for the observational patterns, with which the assimilation data is masked during the training process can be changed. Thus, it is possible for the network to learn Argo-era assimilation patterns while also learning to adjust to pre-Argo era observational density. \\
	The remaining parameters, such as learning rate, batch size and image size, are adopted from the original network set up by \cite{kadow_artificial_2020} and can be viewed with the remainder of the code under \url{https://github.com/slenta/Master-Arbeit}.
	
	\subsection{Data Preprocessing}
	
	The U-net described in section \ref{method:unet} takes three input images in the form of a ground truth assimilation, a binary mask and a masked out assimilation image formatted as an hdf5 file. In order feed the NetCDF4 files containing the assimilation data into the network, they have to be preprocessed to fit the given format. At first, the NA region is cut out of the global assimilation runs. This leaves an irregular 107x124 grid of the NA region from -65 -- -5 \dg W and 20 -- 69 \dg W. For the neural network to be able to process this data, rows of zeros are added until a quadratic 128x128 grid of the NA is reached. Then, the assimilation's temperature profiles are converted to anomalies using the assimilation's average monthly climatology. For this calculation only the time period from 2004 -- 2020 is considered, as this work assumes that the assimilation possesses significantly less uncertainty during the well observed Argo era. Then, EN4 observations from a selected input time are converted to binary masks and multiplied with the data assimilation model data. All three, the masks, the masked data as well as the data assimilations are then stored in an hdf5 file. From there, they can be loaded, split up into batches and fed into the network using a class with the \textit{PyTorch} dataset module.\\
	After the training process, when reconstructing a dataset with monthly observations for a certain timespan, the U-net returns outputs with fully reconstructed monthly temperature profiles.
	
	\subsection{Analysis Approach}
	
	 After the network has been fully trained, visible at the progression of its training curve, it can be given masked temperature profiles for reconstruction. The network will not be trained with all 16 ensemble members. Instead, 4 ensemble members will be held back for later evaluation. This process helps to ensure that the network does not commit to overtraining on its training data. At first, the network will reconstruct the Argo era masked assimilations from one of the held back ensemble members. Each month of assimilation data will be masked with the corresponding observational pattern from that time. The resulting reconstructed OHC can then be summed up over relevant areas and visualized in the form of a timeseries. While the network is trained with the masked assimilations from the entire NA, the data for all other analysis steps is limited to just the SPG. Thus, the neural network is given the opportunity to also learn the SPG surrounding spatial patterns, while the remaining analysis is focused on a region that is of higher interest to current research. Therefore, the results can also be better compared to other OHC estimates. Given the high observational density during the Argo era, this reconstruction can be considered the first step for the neural network towards the reproduction of the assimilation OHC estimate. This approach can then be extended to the pre-Argo era, if the neural network is able to reproduce the assimilation ensemble within its uncertainty range during the Argo era. Analyzing the reconstructed OHC timeseries as well as maps from individual timesteps, the machine learning reconstructions can also be used to discuss inconsistencies and uncertainties within the assimilation. Differences between both OHC estimates can then be taken as indicator for uncertainties that are not mirrored within the assimilation's ensemble spread. Finally, I will examine, whether the neural network is able to directly reconstruct observations in the pre-Argo and Argo era as well as more recent observations it has not trained on within the assimilation's uncertainty range. The results will then be evaluated in regard to the neural network's potential to reproduce or replace EnKF assimilations. Additionally, its implications on a new uncertainty estimate on the EnKF assimilations OHC outside of its ensemble spread will be considered. 

	
	
	
\end{document}